% Proposta da revisionare, NON inclusa automaticamente nei capitoli.
% Richiede la voce selvaraju2017gradcam di heat_map/references.bib.
\subsection{Visualizzazione qualitativa dei contributi alla classificazione}

Per affiancare all'analisi quantitativa una lettura spaziale delle
predizioni viene utilizzato Grad-CAM \cite{selvaraju2017gradcam}.
Il metodo costruisce una mappa di contributi positivi al logit
di una classe, utilizzando le attivazioni dell'ultimo blocco
convoluzionale e il gradiente del punteggio rispetto a esse.
Indicando con \(A^k_{ij}\) l'attivazione spaziale del canale \(k\)
e con \(z_c\) il logit della classe \(c\), si calcolano
\[
 \alpha_k^c=\frac{1}{HW}\sum_{i,j}
       \frac{\partial z_c}{\partial A^k_{ij}},
 \qquad
 L^c_{ij}=\operatorname{ReLU}\left(\sum_k\alpha_k^c A^k_{ij}\right).
\]
La rappresentazione utilizzata è l'output di \texttt{layer4}
della ResNet18 congelata, con 512 canali e risoluzione \(7\times7\).
La mappa viene interpolata sul crop di test \(224\times224\);
l'interpolazione non aumenta la risoluzione informativa della mappa.

B2, Q e QF ricevono la stessa rappresentazione dell'immagine.
Il gradiente viene calcolato attraverso la testa finale del singolo
modello, includendo il circuito in Q e QF. Per rendere il confronto
coerente, tutti i pannelli di una stessa immagine spiegano il logit
della classe vera, anche quando una predizione è errata.
Non vengono aggiornati pesi né ricalcolati split di training.

Prima della visualizzazione vengono verificati gli hash delle
immagini, l'identità dei checkpoint e l'accordo tra le predizioni
ricostruite e quelle archiviate. Gli esempi ordinari sono selezionati
con una regola deterministica per classe; i casi aggiuntivi scelti
per disaccordo tra modelli sono identificati separatamente.
Le mappe hanno finalità illustrativa: non sono attention map
native, non costituiscono una prova causale e non consentono
di dedurre una superiorità generale dell'architettura da pochi esempi.
